\section*{Hoofdstuk \ref{ch:g11:redox} --- Oxidatie en reductie}

\begin{solution}{exo:g11:redox:1}
\omterm{def:g11:redox:oxidant}{Oxidator}: een deeltje dat \omterm{def:g9:inside-the-atom:nucleus}{elektronen} kan opnemen. \omterm{def:g11:redox:oxidant}{Reductor}: een deeltje
dat \omterm{def:g9:inside-the-atom:nucleus}{elektronen} kan afstaan. \omterm{def:g11:redox:oxidation}{Oxidatie}: het afstaan van \omterm{def:g9:inside-the-atom:nucleus}{elektronen}.
\omterm{def:g11:redox:oxidation}{Reductie}: het opnemen van \omterm{def:g9:inside-the-atom:nucleus}{elektronen}.
\end{solution}

\begin{solution}{exo:g11:redox:2}
\ce{Fe^{2+}}/\ce{Fe}; \ce{Ag+}/\ce{Ag}; \ce{I2}/\ce{I-};
\ce{Fe^{3+}}/\ce{Fe^{2+}}.
\end{solution}

\begin{solution}{exo:g11:redox:3}
\ce{Ag+ + e- <=> Ag}; \ce{Al^{3+} + 3e- <=> Al};
\ce{Fe^{3+} + e- <=> Fe^{2+}}; \ce{Cl2 + 2e- <=> 2Cl-}.
\end{solution}

\begin{solution}{exo:g11:redox:4}
Een \omterm{def:g11:redox:oxidation}{oxidatie} (er worden \omterm{def:g9:inside-the-atom:nucleus}{elektronen} afgestaan); ijzer werkt als \omterm{def:g11:redox:oxidant}{reductor}.
\end{solution}

\begin{solution}{exo:g11:redox:5}
Zink wordt geoxideerd en is de \omterm{def:g11:redox:oxidant}{reductor}; \ce{Cu^{2+}} wordt gereduceerd
en is de \omterm{def:g11:redox:oxidant}{oxidator}. Twee \omterm{def:g9:inside-the-atom:nucleus}{elektronen} per zinkatoom.
\end{solution}

\begin{solution}{exo:g11:redox:6}
Begin bij \ce{H2O2 -> H2O}: % ce-unbalanced-ok (step 1 of the method)
zuurstof: twee O links, dus \ce{2H2O} rechts; waterstof: 2 links, 4
rechts, dus \ce{2H+} links; lading: $+2$ links, 0 rechts, dus twee
\omterm{def:g9:inside-the-atom:nucleus}{elektronen} links:
\ce{H2O2 + 2H+ + 2e- <=> 2H2O}.
\end{solution}

\begin{solution}{exo:g11:redox:7}
\ce{I2 + 2e- -> 2I-} en \ce{2S2O3^{2-} -> S4O6^{2-} + 2e-}; elk twee
\omterm{def:g9:inside-the-atom:nucleus}{elektronen}, dus \ce{I2 + 2S2O3^{2-} -> 2I- + S4O6^{2-}}.
\end{solution}

\begin{solution}{exo:g11:redox:8}
\ce{Cu -> Cu^{2+} + 2e-} en $2\times$ (\ce{Ag+ + e- -> Ag}):
\ce{Cu + 2Ag+ -> Cu^{2+} + 2Ag}. De \ce{Cu^{2+}}-ionen die ontstaan, zijn
blauw.
\end{solution}

\begin{solution}{exo:g11:redox:9}
\ce{C6H8O6 -> C6H6O6 + 2H+ + 2e-} en \ce{I2 + 2e- -> 2I-}:
\ce{C6H8O6 + I2 -> C6H6O6 + 2I- + 2H+}. Vitamine C is de \omterm{def:g11:redox:oxidant}{reductor}.
\end{solution}

\begin{solution}{exo:g11:redox:10}
\ce{Al -> Al^{3+} + 3e-} geeft 3 \omterm{def:g9:inside-the-atom:nucleus}{elektronen}, \ce{Cu^{2+} + 2e- -> Cu}
neemt er 2 op; het kleinste gemeenschappelijke veelvoud is 6, dus de
eerste wordt met 2 vermenigvuldigd en de tweede met 3:
\ce{2Al + 3Cu^{2+} -> 2Al^{3+} + 3Cu}.
\end{solution}

\begin{solution}{exo:g11:redox:11}
\ce{H2O2 + 2H+ + 2e- -> 2H2O} en $2\times$ (\ce{Fe^{2+} -> Fe^{3+} + e-}):
\ce{H2O2 + 2H+ + 2Fe^{2+} -> 2H2O + 2Fe^{3+}}.
\end{solution}

\begin{solution}{exo:g11:redox:12}
$n(\ce{Zn}) = 1.30/65.4 = \qty{1.99e-2}{mol}$. Per geoxideerd zinkatoom
slaat er één koperatoom neer, dus $n(\ce{Cu}) = \qty{1.99e-2}{mol}$
en $m(\ce{Cu}) = 1.99\times 10^{-2} \times 63.5 = \qty{1.26}{g}$.
\end{solution}

\begin{solution}{exo:g11:redox:13}
\ce{Cr2O7^{2-} + 14H+ + 6e- -> 2Cr^{3+} + 7H2O} en
$6\times$ (\ce{Fe^{2+} -> Fe^{3+} + e-}):
\ce{Cr2O7^{2-} + 14H+ + 6Fe^{2+} -> 2Cr^{3+} + 6Fe^{3+} + 7H2O}.
Zes \ce{Fe^{2+}} per dichromaation: $6 \times 0.010 = \qty{0.060}{mol}$.
\end{solution}

\begin{solution}{exo:g11:redox:14}
\ce{2ClO- + 4H+ + 2e- -> Cl2 + 2H2O} en \ce{2Cl- -> Cl2 + 2e-}; optellen
en door 2 delen: \ce{ClO- + Cl- + 2H+ -> Cl2 + H2O}. Een zuur levert de
\ce{H+}-ionen, en bij de reactie komt chloor vrij, een giftig gas:
vandaar de waarschuwing dat je de twee producten nooit mag \omterm{def:g2:mixing-and-dissolving:mix}{mengen}.
\end{solution}

\begin{solution}{exo:g11:redox:15}
\omterm{def:g9:inside-the-atom:nucleus}{Elektronen} kunnen niet vrij door de oplossing reizen
(\cref{prop:g11:redox:no-free-electrons}): de \omterm{def:g9:inside-the-atom:nucleus}{elektronen} die de
zinkatomen afstaan, kunnen alleen opgenomen worden door een
\ce{Cu^{2+}}-ion dat het \omterm{def:g9:periodic-table-first-look:metal}{metaal} raakt. Het koperatoom ontstaat dus aan het
oppervlak van de strip, en blijft daar.
\end{solution}

\begin{solution}{pb:g11:redox:1}
\textbf{1.} \ce{Cr2O7^{2-}}/\ce{Cr^{3+}} en \ce{CH3COOH}/\ce{CH3CH2OH}.

\textbf{2.} Ethanol wordt geoxideerd: het is de \omterm{def:g11:redox:oxidant}{reductor}.

\textbf{3.} \ce{Cr2O7^{2-} + 14H+ + 6e- <=> 2Cr^{3+} + 7H2O}.

\textbf{4.} Koolstof klopt (twee aan elke kant). Zuurstof: twee links,
één rechts, dus \ce{H2O} rechts. Waterstof: vier links, $6 + 2 = 8$
rechts, dus \ce{4H+} links. Lading: $+4$ links, 0 rechts, dus vier
\omterm{def:g9:inside-the-atom:nucleus}{elektronen} links:
\ce{CH3COOH + 4H+ + 4e- <=> CH3CH2OH + H2O}.

\textbf{5.} Twaalf \omterm{def:g9:inside-the-atom:nucleus}{elektronen}: $2\times$ de eerste, $3\times$ de tweede
omgekeerd:
\ce{2Cr2O7^{2-} + 3CH3CH2OH + 16H+ -> 4Cr^{3+} + 3CH3COOH + 11H2O}.

\textbf{6.} Oranje \ce{Cr2O7^{2-}}-ionen worden gereduceerd tot groene
\ce{Cr^{3+}}-ionen, overal waar de ethanol ze bereikt.

\textbf{7.} $M = 2\times 39.1 + 2\times 52.0 + 7\times 16.0 =
\qty{294.2}{g/mol}$.

\textbf{8.} $n = 2.0\times 10^{-3}/294.2 = \qty{6.8e-6}{mol}$.

\textbf{9.} Drie ethanolmoleculen op twee dichromaationen:
$n = \tfrac{3}{2}\times 6.8\times 10^{-6} = \qty{1.0e-5}{mol}$.

\textbf{10.} $M(\ce{C2H6O}) = \qty{46.0}{g/mol}$, dus
$m = 1.02\times 10^{-5}\times 46.0 = \qty{4.7e-4}{g}$, ongeveer
\qty{0.47}{mg}.

\textbf{11.} $0.25/0.47 \approx 0.53$: ongeveer \qty{53}{\percent} van
het dichromaat.

\textbf{12.} $0.53 \times 10 \approx \qty{5.3}{cm}$.

\textbf{13.} $\qty{0.25}{mg} \times 2100 = \qty{525}{mg}$ per liter
bloed, ongeveer \qty{0.53}{g/L}.

\textbf{14.} De adem komt eerst langs de korrels bij de inlaat; daar is
het dichromaat in overmaat en haalt het alle ethanol weg, zodat de
reactie in een eerste stuk van het buisje volledig is, dat helemaal groen
wordt, voordat er ethanol verder komt.

\textbf{15.} De reactie is aflopend: de groene \ce{Cr^{3+}}-ionen worden
niet weer dichromaat, dus een gebruikt buisje kan niet opnieuw meten.

% ledger: ghs:K2Cr2O7 (H350 among its hazard statements)
\textbf{16.} Het is acuut giftig (GHS06) en ernstig gevaarlijk voor de
gezondheid (GHS08: tot de gevarenaanduidingen behoort ``kan kanker
veroorzaken''), en daarnaast bijtend en giftig voor in het water levende
organismen; een apparaat dat door het publiek gebruikt wordt, mag het
niet bevatten.

\textbf{17.} \ce{O2 + 4H+ + 4e- -> 2H2O} en
\ce{CH3CH2OH + H2O -> CH3COOH + 4H+ + 4e-}:
\ce{CH3CH2OH + O2 -> CH3COOH + H2O}.

\textbf{18.} Vier \omterm{def:g9:inside-the-atom:nucleus}{elektronen}.

\textbf{19.} $n = 0.25\times 10^{-3}/46.0 = \qty{5.4e-6}{mol}$
ethanol; \omterm{def:g9:inside-the-atom:nucleus}{elektronen} $4n = \qty{2.2e-5}{mol}$, dat is
$2.2\times 10^{-5}\times 6.02\times 10^{23} = \num{1.3e19}$ \omterm{def:g9:inside-the-atom:nucleus}{elektronen}.

\textbf{20.} $Q = 1.3\times 10^{19}\times 1.60\times 10^{-19} \approx
\qty{2.1}{C}$. Elk ethanolmolecuul stuurt precies vier \omterm{def:g9:inside-the-atom:nucleus}{elektronen} door
de draad, dus de lading is evenredig met de hoeveelheid ethanol: wie de
lading meet, meet de ethanol.
\end{solution}
